\chapter{The Collaboration: A Retrospective}

\section{The convergence}

Two independent projects started from opposite ends of the software stack and
built toward the same shape, without coordination.

\paragraph{sitas} began in userspace, asking: \emph{what does shared-nothing,
shard-per-core concurrency look like when it is built out of Rust ownership
instead of locks?} Its answer was a discipline --- each shard owns its state,
only the owning shard mutates it, everything crossing a shard boundary is an
owned value moved through an explicit typed message. On top of that it grew a
custom async executor, per-shard reactors, \texttt{io\_uring} completion I/O,
and snapshot-based observability.

\paragraph{CharlotteOS} began in the kernel, asking: \emph{hardware is
asynchronous with respect to the CPU, so why is the OS synchronous?} Its answer
was to make the kernel async-first --- cheap threads, an observer/waker event
model, and system calls that are asynchronous by default, returning a
completion capability that userspace can wait on.

Both projects independently converged on the same primitive: a waitable
completion handle. sitas called it \texttt{Reply<T>} / \texttt{Notify};
CharlotteOS called it \texttt{Observer} / \texttt{Waker}. They are the same
object with different names.

The insight that drove the collaboration: \textbf{a userspace runtime whose
async model agrees with the kernel it runs on removes the ``retrofit tax''}
--- the async-signal-safety, thread pools, and \texttt{io\_uring}-over-blocking
emulation that exist only because the runtime and the OS disagree about what is
fundamental.

\section{The structural correspondence}

The reason the projects belong together is deeper than analogy: the pieces
already existed on both sides, and they lined up almost one-to-one.

\begin{table}[h]
\centering
\begin{tabularx}{\textwidth}{l|l|l}
\textbf{sitas concept} & \textbf{Role} & \textbf{CharlotteOS equivalent} \\
\hline
Shard = OS thread pinned to core & Unit of isolation & Logical Processor (LP) + \texttt{PerLp<T>} \\
\texttt{ShardId} explicit placement & Addressing & \texttt{LpId} \\
Shard mailbox (\texttt{ShardSender}/\texttt{ShardReceiver}) & Cross-shard transport & \texttt{IPI\_CMD\_QUEUES} (bounded per-LP) \\
Custom executor \texttt{Waker} re-enqueue & Wake a task & \texttt{Waker(ThreadId)} + \texttt{submit\_ready\_thread} \\
\texttt{Reply<T>} / \texttt{Notify} awaitable completion & Await a result & \texttt{Observable} / \texttt{Observer} + \texttt{TimerEvent} \\
Reactor idle wait & Block until events & GIC/APIC IRQ dispatch $\rightarrow$ \texttt{cond\_yield\_lp} \\
\texttt{io\_uring} completion I/O & Async syscalls & Completion-capability ABI (design goal) \\
\texttt{ShardLocal<T>} (no mutex, owner-checked) & Shard-owned mutable state & \texttt{ShardLocal<T>} in-kernel (sitas-derived) \\
CPU pinning & Bind shard to core & LP \emph{is} a core (not advisory) \\
\end{tabularx}
\end{table}

\texttt{sitas} was, in effect, re-deriving in userspace the runtime that
CharlotteOS was trying to be natively. Where sitas fought to build a reactor on
top of a synchronous OS, CharlotteOS offered the reactor as the substrate.

\section{The exploration phases (A $\rightarrow$ C $\rightarrow$ B)}

The collaboration was structured as three time-boxed spikes, pursued in
dependency order, each building on the prior. Nothing was merged until it
earned its place. All spikes were verified on QEMU (AArch64 and, later, x86-64).

\subsection{Phase A --- OS backend seam (sitas)}

\textbf{Goal:} extract the reactor's OS dependency into an explicit trait, so
the executor can be backed by anything that can block until events and wake
from elsewhere --- not just Unix \texttt{epoll}/\texttt{kqueue}.

\textbf{Deliverable:} the \texttt{ReactorBackend} trait (in sitas's
\texttt{reactor\_backend.rs}). Three methods capture everything the executor
needs from the OS:

\begin{itemize}
    \item \texttt{waker()} --- obtain a cloneable waker.
    \item \texttt{wait(read, write, timeout)} $\rightarrow$ \texttt{Event} ---
        block until an interest becomes ready, the reactor is woken, or a
        deadline elapses.
    \item \texttt{wake()} --- wake a blocked reactor from another thread/core.
\end{itemize}

A mock in-memory backend proved a non-Unix backend is implementable behind the
same interface. The existing Unix \texttt{OsReactor} / \texttt{OsWaker} /
\texttt{OsEvent} satisfied the contract with zero changes.

\textbf{Key finding:} abstracting ``block until an event'' was trivial.
Abstracting ``\emph{what is a waitable interest}'' was the real work. On Unix
the interest is a \texttt{RawFd}; the trait's \texttt{Handle} associated type
is where the design pressure concentrates. For CharlotteOS, \texttt{Handle} is a
completion capability.

\subsection{Phase C --- Async syscall ABI (CharlotteOS kernel)}

\textbf{Goal:} design the userspace$\leftrightarrow$kernel boundary for
asynchronous system calls, using sitas as the executable specification. If
sitas's shard model, completion futures, and cross-shard submit map cleanly
onto the ABI, the ABI is right.

\textbf{Deliverable:} five kernel-internal operations that later became the
syscall table:

\begin{enumerate}
    \item \texttt{submit(op, args, buffers)} $\rightarrow$ \texttt{CompletionCap}
        (a per-address-space handle naming in-flight work).
    \item \texttt{wait(cq, min\_complete, deadline)} $\rightarrow$ count (the
        reactor's only sleep).
    \item \texttt{wake(cq)} (cross-shard wake via IPI).
    \item \texttt{cancel(cap)} (deferred buffer reclaim).
    \item \texttt{close(cap)} (release a completed slot).
\end{enumerate}

Each built on facilities that already existed in the kernel: the per-AS
\texttt{IdTable} for the capability table, the \texttt{Observable}/\texttt{Observer}
mechanism for completing operations, the per-LP \texttt{RoundRobin} scheduler
for blocking and waking threads, and the bounded \texttt{IPI\_CMD\_QUEUES} for
cross-core communication.

\textbf{Key infrastructure built in this phase:}

\begin{itemize}
    \item \textbf{Completion-capability table} (\texttt{completion/mod.rs}) ---
        per-address-space \texttt{IdTable<Arc<Completion>>}, bounded for
        submission backpressure.
    \item \textbf{CQ ring} (\texttt{completion/cq.rs}) --- a shared-memory
        io\_uring-style ring buffer for zero-syscall completion delivery.
        Single 4 KiB page, header + entry array. The kernel writes entries;
        userspace reads them without syscalls.
    \item \textbf{Syscall dispatch} (\texttt{syscall/mod.rs}) --- the AArch64
        \texttt{sync\_dispatcher} decodes \texttt{ESR\_EL1.EC} and routes SVC
        instructions to a dispatch table. A \texttt{TrapFrame} captures the user
        register state.
    \item \textbf{Bounded IPI queues} --- \texttt{IPI\_CMD\_QUEUES} became
        bounded per-LP \texttt{ConcurrentQueue<256>} with first-class
        backpressure (\texttt{try\_push\_to} returns \texttt{Err(rpc)} on full).
    \item \textbf{\texttt{IpiRpc::Closure}} variant --- arbitrary
        \texttt{FnOnce() + Send} work dispatched to a target LP, the seed of
        \texttt{ShardMailbox<M>}.
\end{itemize}

\subsection{Phase B --- sitas's discipline inside the kernel}

\textbf{Goal:} adopt sitas's shared-nothing invariants as the kernel's own
concurrency doctrine. CharlotteOS already had the ingredients; sitas
contributed the rules.

\textbf{Deliverables:}

\begin{itemize}
    \item \textbf{\texttt{ShardLocal<T>}} --- a lock-free per-LP container
        replacing \texttt{PerLp<T>}'s \texttt{RwLock} with \texttt{UnsafeCell<T>}
        and two runtime assertions: owner-check (the caller must be on the
        owning LP) and re-entrancy guard (borrow flag). Closure-based access
        ensures \texttt{\&mut T} never escapes.
    \item \textbf{\texttt{ShardMailbox<M>}} --- a typed, bounded per-LP queue
        with cloneable \texttt{ShardSender<M>} and single-consumer
        \texttt{ShardReceiver<M>}, built on the bounded IPI infrastructure.
        \texttt{try\_send} returns \texttt{Err(m)} on backpressure.
\end{itemize}

\subsection{Beyond the plan --- what was built additionally}

\begin{itemize}
    \item \textbf{Exit-observer completion} --- the ABI's intended mechanism:
        \texttt{submit\_worker} registers a completion capability as the
        worker thread's exit-observer. The thread exiting \emph{is} the
        completion event. This required fixing the thread lifecycle
        (deferred reaper) so threads could return cleanly.
    \item \textbf{Real-EL0 user thread} --- the first userspace thread running
        at EL0, invoking \texttt{SVC} from compiled Rust code, with the kernel
        dispatching the syscall and returning.
    \item \textbf{x86-64 parity} --- serial console (16550 UART), SYSCALL entry
        trampoline, unstubbed IPI handlers, LA57 paging-mode fix. The same
        self-test suite passes on both architectures.
    \item \textbf{16550 UART serial driver} --- x86-64 headless boot was
        completely silent before this; debugging was impossible.
    \item \textbf{Multi-LP bring-up} --- kernel boots with \texttt{-smp 2} on
        AArch64 (cross-core SGI delivery is blocked by a QEMU GICv3 emulation
        defect; self-targeted SGIs and per-core timers work on all LPs).
    \item \textbf{sitas-core no\_std refactor} --- the sitas executor, reactor
        traits, ShardLocal, ShardMailbox, and completion primitives were split
        from the Unix backend into a \texttt{no\_std + alloc} crate. A lock-free
        \texttt{ringbuf} replaced \texttt{concurrent-queue}. The
        \texttt{ShardRuntime} trait abstracts thread spawning.
    \item \textbf{sitas-charlotte backend} --- implements
        \texttt{ReactorBackend} and \texttt{ShardRuntime} against the kernel's
        async syscall ABI. The user binary links sitas-core + sitas-charlotte
        and runs \texttt{ShardedExecutor} on CharlotteOS.
\end{itemize}

\section{The async syscall ABI in detail}

The ABI is the concrete interface where both projects meet. It is specified
against sitas's \texttt{ReactorBackend} trait and implemented in the
CharlotteOS kernel.

\subsection{The five operations}

\begin{description}
    \item[\texttt{submit(asid, op, buffer)} $\rightarrow$ \texttt{CompletionCap}]
        Returns immediately with a capability naming the in-flight operation.
        Ownership of the buffer transfers to the kernel until a terminal
        completion posts it back. Returns \texttt{SubmitError::WouldBlock} on
        submission backpressure (capability table full).

    \item[\texttt{wait(asid, cap)} $\rightarrow$ \texttt{Result<(), CapError>}]
        Blocks the calling thread until the specified capability reaches
        terminal completion. Internally registers the thread's \texttt{Waker}
        as an \texttt{Observer} of the \texttt{Completion} (which is
        \texttt{Observable}) and yields.

    \item[\texttt{poll(asid, cap)} $\rightarrow$ \texttt{Option<Completed>}]
        Non-blocking check. Returns \texttt{None} while in flight; returns
        \texttt{Some(Completed \{ result, buffer \})} when terminal. Buffer
        ownership returns to the caller.

    \item[\texttt{cancel(asid, cap)} $\rightarrow$ \texttt{CancelState}]
        Requests cancellation. If in flight, returns \texttt{CancelRequested}:
        the kernel retains any transferred buffer until it posts a terminal
        \texttt{Cancelled} completion (deferred reclaim). If already complete,
        returns \texttt{AlreadyComplete}.

    \item[\texttt{close(asid, cap)} $\rightarrow$ \texttt{Result<(), CapError>}]
        Frees the capability-table slot after the terminal completion was
        posted or the result was consumed by \texttt{poll}. Fails with
        \texttt{NotComplete} for in-flight operations.
\end{description}

\subsection{The exit-observer completion path}

The ABI's intended completion mechanism is declarative: \textbf{the worker
thread exiting IS the completion event.} \texttt{submit\_worker(asid,
worker\_entry, result)} spawns a worker thread and registers the returned
capability as the worker's exit-observer. The worker performs its work and
returns; the thread's \texttt{Thread::Drop} fires the exit-observer, which
calls \texttt{complete()} and wakes any waiter. The worker never references the
capability explicitly.

This matches the design comment in
\texttt{scheduler/threads/mod.rs:63--69} and is implemented as
\texttt{CompletionExitObserver} in
\texttt{crates/catten/src/completion/mod.rs}.

\subsection{The CQ ring}

A shared-memory io\_uring-style ring buffer for zero-syscall completion
delivery. A single 4 KiB page contains a header (head, tail, capacity,
overflow) and a circular array of entries. The kernel (producer) writes entries
via \texttt{write(cap, result)}; userspace (consumer) drains via
\texttt{read() $\rightarrow$ Option<CqEntry>} -- zero syscalls.

\subsection{Cancellation and the buffer-reclaim contract}

When userspace drops a \texttt{CompletionCap} without awaiting it:

\begin{enumerate}
    \item The kernel marks the operation cancelling but \textbf{retains the
        buffer} --- it may still be the target of DMA or kernel access.
    \item A terminal \texttt{Cancelled} completion is eventually posted,
        handing the buffer back.
    \item On teardown, outstanding buffers are leaked rather than freed while
        the kernel may touch them.
\end{enumerate}

This is the kernel-side mirror of sitas's \texttt{defer\_buffer\_drop} and
\texttt{mem::forget}-on-timeout rule.

\subsection{Backpressure}

\begin{itemize}
    \item \textbf{Submission:} the per-AS capability table is bounded.
        \texttt{submit} returns \texttt{WouldBlock} when full.
    \item \textbf{Completion:} the CQ ring is bounded. When full, the kernel
        does \emph{not} drop completions; it stops producing and marks the
        queue overflow-pending.
    \item \textbf{Cross-LP:} \texttt{IPI\_CMD\_QUEUES} are bounded per target
        LP. \texttt{try\_push\_to} returns \texttt{Err(rpc)} on full.
\end{itemize}

\section{The bugs we found (and the real value of running on hardware)}

Bootstrapping an async-first kernel on real hardware (QEMU) surfaced a series
of latent bugs that static analysis and self-tests never caught. Each fix
enabled the next layer of work. The following are the most instructive, in the
order they were discovered.

\subsection{Panic handler recursed through heap-dependent \texttt{logln}}

\textbf{Symptom:} any fault during early boot (before the heap allocator was
ready) silently triple-faulted instead of printing a panic message --- making
\emph{every} early bug look like a mysterious hang.

\textbf{Root cause:} the panic handler used \texttt{logln!}, which routes
through \texttt{INT\_STATE} --- a \texttt{LazyLock} whose first-use
initialization allocates from the kernel heap (\texttt{alloc::vec!}). A fault
before the heap was ready caused the panic handler's \texttt{logln!} to trigger
a heap allocation $\rightarrow$ another fault $\rightarrow$ re-entered the
panic handler $\rightarrow$ infinite recursion $\rightarrow$ stack overflow
$\rightarrow$ triple fault.

\textbf{Fix:} the panic handler now uses \texttt{early\_logln!} (direct serial
writes, no heap, no \texttt{INT\_STATE}). A panic handler must never depend on
the allocator it might be reporting a failure about.

This was the single most valuable fix of the entire project: it unmasked every
subsequent fault as a visible panic instead of a silent triple-fault.

\subsection{Blocked threads lost their execution context}

\textbf{Symptom:} a thread that blocked in \texttt{wait()} (or \texttt{sleep()})
restarted from its entry point when woken, instead of resuming where it
blocked.

\textbf{Root cause:} \texttt{block\_thread} cleared the LP scheduler's
\texttt{current\_handle} for the running thread, so the subsequent
\texttt{cond\_yield\_lp} read \texttt{curr\_tid == None} and
\texttt{switch\_ctx(null, ...)} \textbf{skipped saving} the blocking thread's
execution context.

\textbf{Fix:} a self-blocking thread now stays \texttt{current\_handle} until
the yield saves its context; \texttt{RoundRobin::next} declines to re-queue a
\texttt{Blocked} thread. This was a foundational fix: \texttt{wait()} and
\texttt{sleep()} had never been exercised before.

\subsection{RwLock held-across-call deadlocks in \texttt{sleep()} and
\texttt{abort()}}

\textbf{Symptom:} \texttt{sleep()} worked reliably with diagnostic logging but
hung without it --- a classic heisenbug.

\textbf{Root cause:} \texttt{sleep()} held a \texttt{SYSTEM\_SCHEDULER.read()}
guard (plus the LP scheduler lock) across
\texttt{SYSTEM\_SCHEDULER.write().block\_thread(...)} --- a read-then-write
RwLock self-deadlock. The diagnostic version happened to bind the thread ID
first, releasing the guard. \texttt{abort()} had the same pattern with the LP
scheduler lock.

\textbf{Fix:} bind the value out of the scrutinee first so the temporary
guards drop before the re-locking call.

\subsection{Quantum timer grew the timer queue without bound}

\textbf{Symptom:} after extended operation, the per-LP timer queue contained
thousands of quantum events, consuming unbounded memory.

\textbf{Root cause:} \texttt{clear\_ctx\_switch\_pending} enqueued a new
quantum \texttt{TimerEvent} on every yield, not just on quantum expiry. Manual
yields accumulated quantum events forever.

\textbf{Fix:} an \texttt{armed} flag keeps exactly one quantum event in flight
(the quantum's own firing clears it). Also, \texttt{cond\_yield\_lp} re-arms even
when there is nothing to switch to (previously the timer stopped when only one
thread was runnable), and \texttt{process\_events} stops the level-triggered
timer when the queue drains.

\subsection{A thread cannot free its own kernel stack}

\textbf{Symptom:} a returning kernel thread hung during teardown.

\textbf{Root cause:} \texttt{ThreadContext::Drop} calls
\texttt{deallocate\_stack}, which frees the kernel stack --- but the returning
thread is still executing on that stack when \texttt{abort} drops it.

\textbf{Fix:} a \textbf{deferred reaper}. \texttt{abort} moves the dying thread
to \texttt{DEAD\_THREADS} (\texttt{IdTable::take\_element} extracts it without
dropping), and \texttt{cond\_yield\_lp} calls \texttt{reap\_dead\_threads()}
\emph{after} switching away, so the stack is freed from a different thread's
context.

\subsection{Stack allocator never registered guard pages}

\textbf{Symptom:} \texttt{deallocate\_stack} always failed with
\texttt{InvalidStack}, causing the reaper to leak stacks.

\textbf{Root cause:} \texttt{allocate\_stack} never populated
\texttt{KERNEL\_GUARD\_PAGE\_SET}, so \texttt{validate\_stack} couldn't find
the guards.

\textbf{Fix:} register lower and upper guard pages on allocation, and rewrite
\texttt{deallocate\_stack} to derive the page count from them. Reference-count
guard pages (using \texttt{BTreeMap<VAddr, usize>}) so adjacent stacks sharing
a guard survive until both are freed.

\subsection{LA57 paging-mode mismatch on x86-64}

\textbf{Symptom:} a \texttt{\#GP} with \texttt{RAX = 0xff90000000000000}
(clearly non-canonical under 4-level paging) during kernel-heap large-page
mapping.

\textbf{Root cause:} \texttt{CpuInfo::get\_vaddr\_sig\_bits()} read the CPU's
\emph{maximum} linear address width from CPUID \texttt{0x80000008}
EAX[15:8] -- with \texttt{-cpu max} that is 57 (the CPU supports 5-level
paging). The kernel then selected \texttt{LA\_MAP\_57BIT}, whose
\texttt{kernel\_allocator\_arena} base is \texttt{0xff90000000000000} --
canonical only under 5-level paging. But Limine boots with 4-level paging
(48-bit VAs), under which that address is non-canonical.

\textbf{Fix:} derive the address width from \texttt{CR4.LA57} (bit 12) ---
the \emph{active} paging mode, not the CPU's capability.

\subsection{ASID-0 collision (``the TLB bug that wasn't'')}

\textbf{Symptom:} the first user thread faulted at its entry point with
an instruction abort (\texttt{EC=0x21}, ``no change in exception level'').

\textbf{Root cause:} the global \texttt{ADDRESS\_SPACE\_TABLE} started empty,
so the first user address space got id 0, which equals \texttt{KERNEL\_ASID}.
\texttt{Thread::new} then built a \emph{kernel} thread context, running the
user code at EL1 where \texttt{PXN=1} forbids execution. The EC decode
(\texttt{0x21} = fault without EL change) was the tell --- the code was
never at EL0.

\textbf{Fix:} pre-populate the table with the kernel AS at id 0, so user
spaces always get non-zero ids.

A second bug in the same test: \texttt{sync\_dispatcher} added 4 to
\texttt{ELR\_EL1} on SVC, but AArch64 already sets ELR to the instruction
\emph{after} SVC. The +4 overshot the loop instruction, landing in zeroed
memory (\texttt{EC=0} undefined instruction).

\subsection{x86-64 \texttt{kernel\_thread\_trampoline} halted on thread return}

\textbf{Symptom:} on x86-64, the async demo's worker thread completed its
work but the coordinator never resumed.

\textbf{Root cause:} the x86-64 \texttt{kernel\_thread\_trampoline} entered a
\texttt{hlt} loop when a thread's entry returned, instead of calling
\texttt{abort()} like the AArch64 trampoline. The CPU halted, freezing the LP.

\textbf{Fix:} make the x86-64 trampoline call \texttt{abort} on entry return,
matching AArch64's behavior. (The AArch64 trampoline had always called
\texttt{abort} via \texttt{bl abort} after the entry point returned.)

\section{Verification}

All verification is via QEMU emulation (AArch64 \texttt{virt} and x86-64
\texttt{q35} machines, TCG accelerator). The kernel has no host-side test
framework; all tests are boot-time \textbf{kernel self-tests} that run from
\texttt{main.rs:104} via \texttt{self\_test::run\_self\_tests()}.

\subsection{Self-test suite (both architectures pass)}

\begin{itemize}
    \item \texttt{self\_test/completion.rs} --- buffer ownership, cancel /
        deferred-reclaim, backpressure, close rejection
    \item \texttt{self\_test/syscall.rs} --- dispatch table routing via
        synthetic TrapFrame (all 7 syscalls routed without panicking)
    \item \texttt{self\_test/ipi.rs} --- bounded queue, backpressure, closure
        dispatch
    \item \texttt{self\_test/shard.rs} --- \texttt{ShardLocal<T>} owner-check /
        re-entrancy guard, \texttt{ShardMailbox<M>} send/recv
    \item \texttt{self\_test/el0.rs} --- real-EL0 user thread (AArch64 only),
        compiled Rust binary calling SVC with correct dispatch
    \item \texttt{self\_test/cq.rs} --- CQ ring write/drain/overflow detection
    \item \texttt{self\_test/cq\_completion.rs} --- CQ ring integrated with
        completion subsystem
\end{itemize}

\subsection{Async-syscall demo (both architectures)}

The \texttt{demo.rs} spawns a coordinator + worker from \texttt{bsp\_main}.
The coordinator calls \texttt{submit\_worker}, which spawns a worker and
registers the capability as the worker's exit-observer. The worker performs
async work (\texttt{sleep 50ms}) and returns; the thread-exit fires the
observer, completing the capability and waking the coordinator. Verified
output:

\begin{verbatim}
Testing Complete. All Tests Passed!
[async-demo] coordinator: submitted worker, now awaiting completion...
[async-demo] worker: performing async work (sleep 50ms)
[async-demo] worker: work finished, exiting (completion fires on exit)
[async-demo] coordinator: COMPLETION RECEIVED, result Ok(42)
[async-demo] SUCCESS: async syscall round-trip complete (via thread-exit observer)
\end{verbatim}

\subsection{sitas integration}

The user binary (\texttt{crates/catten-user/}) compiles against
\texttt{sitas-core} (no\_std executor) and \texttt{sitas-charlotte}
(ReactorBackend against kernel syscalls). The binary creates a
\texttt{ShardedExecutor} with \texttt{CharlotteReactor} as the
\texttt{ShardRuntime}, spawning a shard thread via SVC \#7
(\texttt{SPAWN\_THREAD}). The spawned thread writes a sentinel to the shared
result page. Binary size: 144 bytes (up from 48 without sitas). Zero panics.

\subsection{CI}

GitHub Actions (\texttt{.github/workflows/ci.yml}) builds both architectures
with \texttt{-D warnings}, and runs an automated QEMU boot gate on the
\texttt{dev} branch (60-second timeout, grepping for ``Testing Complete. All
Tests Passed!'').

\section{What we learned}

\begin{enumerate}
    \item \textbf{The block-yield-wake cycle is the hard problem.} Designing
        the completion ABI was straightforward compared to making
        \texttt{wait()} correctly save and restore thread state through the
        cooperative scheduler. Every scheduling primitive (\texttt{sleep},
        \texttt{wait}, \texttt{abort}) exercised a new path and surfaced a new
        bug.

    \item \textbf{Booting on real hardware finds bugs that static analysis
        cannot.} The panic-handler recursion, ASID-0 collision, and LA57
        paging mismatch were all invisible in self-tests and logical reasoning.
        Every one was found by booting and observing the actual fault.

    \item \textbf{The observer/waker pattern is the right primitive.} Once
        \texttt{Completion} was made \texttt{Observable} and \texttt{wait()}
        used the same path as \texttt{sleep()}, the entire completion subsystem
        fell into place. The exit-observer mechanism is the natural endpoint:
        a thread exiting \emph{is} the completion event.

    \item \textbf{The async model genuinely removes the retrofit tax.} On
        CharlotteOS, \texttt{wait(cap)} blocks exactly as \texttt{sleep(50ms)}
        blocks --- both register a \texttt{Waker} on an \texttt{Observable} and
        yield. There is no emulation, no thread pool, no signal-handler
        tightrope. The kernel and the runtime agree about what is fundamental.

    \item \textbf{The sitas traits were essential for co-design.} The
        \texttt{ReactorBackend} trait defined in Phase A became the executable
        specification for the async syscall ABI in Phase C. The sitas-charlotte
        backend validates the ABI shapes before any kernel path exists. This
        bidirectional validation is the core of the collaboration.
\end{enumerate}

\section{Current status (July 2026)}

\textbf{Verified on QEMU AArch64 and x86-64:}

\begin{itemize}
    \item Both architectures boot, all self-tests pass, exit-observer
        async demo succeeds.
    \item Completion subsystem: submit, complete, poll, wait, cancel, close,
        \texttt{submit\_worker}, exit-observer, CQ ring.
    \item Syscall dispatch: AArch64 SVC handler (7 syscalls). x86-64 SYSCALL
        trampoline exists (connected to dispatch table).
    \item IPI: bounded per-LP queues, backpressure, typed closure dispatch.
    \item \texttt{ShardLocal<T>} / \texttt{ShardMailbox<M>} in-kernel.
    \item Scheduler: RoundRobin, quantum timer, block/yield/wake, deferred
        reaper, per-LP thread pinning.
    \item sitas-core no\_std crate + sitas-charlotte ReactorBackend +
        \texttt{ShardRuntime::spawn\_shard} (SVC \#7).
    \item CI: GitHub Actions builds both arches, QEMU boot gate.
\end{itemize}

\textbf{Known limitations (documented, scoped):}

\begin{itemize}
    \item x86-64 ring-3 / SYSCALL userspace (GS base, IRETQ return) --- not
        yet exercised.
    \item Full \texttt{ShardSet}/\texttt{Sharded} threading layer (the
        \texttt{ShardRuntime} trait needs to be threaded through the
        \texttt{ShardSet::start} call-chain; currently \texttt{ShardedExecutor}
        works directly but the higher-level service modules still use
        \texttt{thread::spawn}).
    \item Multi-LP cross-core IPI (QEMU GICv3 SGI delivery defect --- self-IPIs
        and per-core timers work on all LPs).
\end{itemize}

\section{Repositories and branches}

\begin{itemize}
    \item \textbf{sitas} (\texttt{reactor-handle-seam}):
        \texttt{ReactorBackend} trait, ABI reference model
        (\texttt{charlotte\_abi.rs}), \texttt{sitas-core} no\_std crate,
        \texttt{sitas-charlotte} crate.
    \item \textbf{CharlotteOS} (\texttt{dev}): consolidated kernel with
        completion subsystem, syscall dispatch, bounded IPI queues,
        ShardLocal, ShardMailbox, CQ ring, real-EL0 test, exit-observer
        completion, \texttt{SPAWN\_THREAD} syscall, sitas user binary.
\end{itemize}

\endinput
